\chapter{Stravinsky, Igor}
\label{cha:stravinsky}

\href{https://www.kcatalog.org/index.php/browse-chapters/kcatalog}{KCatalogue} of Kirchmeyer.

\section{Opera/Theatre}

\noindent{The Nightingale (Le Rossignol), 3-act opera (1914), K018} \\
\noindent{Renard, a burlesque for 4 pantomimes and chamber orchestra (1916), K023} \\
\noindent{L'Histoire du soldat (The Soldier's Tale), for chamber ensemble and three speakers (1918), K029} \\
\noindent{Mavra, one-act opera (1922), K039} \\
\noindent{Oedipus rex, 2-act opera-oratorio (1927), K047} \\
\noindent{Perséphone, mélodrame for speaker, soloists, chorus, and orchestra (1933), K056} \\
\noindent{The Rake's Progress, 3-act opera (1951), K078} \\
\noindent{The Flood, television opera (1962), K098} \\

\section{Ballet}

\noindent{The Firebird (L'oiseau de feu) (1910; rev. 1919, 1945), K010} \\
\noindent{Petrushka (1911, rev. 1947), K012}  \textit{Orchestre de Paris, Klaus M\"{a}kel\"{a}}\marginnote{01/2001, max} \\
\noindent{The Rite of Spring (Le sacre du printemps) (1913; rev. 1947, 1967), K015} \\
\noindent{Les Noces (The Wedding), for soloists, choir, four pianos and percussion (1914–17 and 1919–23), K040} \\
\noindent{Pulcinella, for chamber orchestra and soloists (1920), K034} \\
\noindent{Apollo (Apollon musagète), for string orchestra (1928; rev. 1947), K048}  \textit{Nieuw Sinfonietta Amsterdam, Lev Markiz}\marginnote{high} \\
\noindent{Le Baiser de la fée (The Fairy's Kiss) (1928; rev. 1950), K049}  \textit{London PO, Vladimir Jurowski}\marginnote{max} \\
\noindent{Jeu de cartes (Card Game) (1936), K059} \\
\noindent{Danses concertantes for chamber orchestra (1941–42), K063} \\
\noindent{Ode (1943), K066} \\
\noindent{Scènes de ballet (1944), K069} \\
\noindent{Orpheus, for chamber orchestra (1947), K076} \\
\noindent{Agon (1957), K088} \\

\section{Orchestral}

\noindent{Symphony in E-flat major, Op. 1 (1907), K003} \\
\noindent{Scherzo fantastique, Op. 3 (1908), K005} \\
\noindent{Feu d'artifice (Fireworks), Op. 4 (1908), K007} \\
\noindent{Funeral Song (Chant funèbre), Op. 5 (1908), K008} \\
\noindent{Le chant du rossignol (Song of the Nightingale) (1917), K026} \\
\noindent{Suite from Pulcinella (1920; rev. 1947), K034} \\
\noindent{Symphonies of Wind Instruments (1920), K036} \\
\noindent{8 Instrumental Miniatures for 15 Players (1963, orchestration of Les cinq doigts), K037} \\
\noindent{Suite No. 2 for chamber orchestra (1921, arrangement of Trois pièces faciles and Cinq pièces faciles No. 5), K038} \\
\noindent{Suite No. 1 for chamber orchestra (1925, arrangement of Cinq pièces faciles Nos. 1–4), K045} \\
\noindent{Divertimento (Suite from Le Baiser de la fée, 1934), K049} \textit{BBC PO, Andrew Davis}\marginnote{max} \\
\noindent{Quatre études for orchestra (1928, arrangement of Three Pieces for String Quartet and Étude pour pianola), K051} \\
\noindent{Concerto in E-flat Dumbarton Oaks for chamber orchestra (1938), K060}  \textit{Nieuw Sinfonietta Amsterdam, Lev Markiz}\marginnote{high} \\
\noindent{Symphony in C (1940), K061} \textit{BBC PO, Andrew Davis}\marginnote{max} \\
\noindent{Tango for chamber orchestra (1953, arrangement of 1940 work for piano), K062} \\
\noindent{Circus Polka (1944), K064} \textit{BBC PO, Andrew Davis}\marginnote{max} \\
\noindent{Four Norwegian Moods (1942), K065} \\
\noindent{Scherzo à la russe for orchestra (1944, also a version for Paul Whiteman's band), K070} \\
\noindent{Symphony in Three Movements (1945), K073} \textit{BBC PO, Andrew Davis}\marginnote{max} \\
\noindent{Concerto in D "Basle" for string orchestra (1946), K075}  \textit{Nieuw Sinfonietta Amsterdam, Lev Markiz}\marginnote{high} \\
\noindent{Greeting Prelude for orchestra (1955), for the 80th birthday of Pierre Monteux, K085}  \textit{BBC PO, Andrew Davis}\marginnote{max} \\
\noindent{Variations (Aldous Huxley in Memoriam) (1963/1964), K103} \\
\noindent{Canon on a Popular Russian Tune (1965), K105} \\

\section{Concertante}
\noindent{Scherzo for Violin and Orchestra (1882), K001} \\
\noindent{Piano Concerto No. 1 in E-flat major, Op. 6 (1911–12), K010} \\
\noindent{Violin Concerto in D major (1931), K047}  \textit{Frank Peter Zimmermann, Bamberger Symphoniker, Jakob H\v{r}usa}\marginnote{02/2024, max} \\
\noindent{Capriccio for Piano and Orchestra (1929), K050} \\
\noindent{Concerto for Two Solo Pianos (1935), K057} \\
\noindent{Ebony Concerto for Clarinet and Jazz Band (1945), K072} \\
\noindent{Concerto in D for Violin and Orchestra (1947), K076} \\
\noindent{Canticum Sacrum ad honorem Sancti Marci Nominis (1955), K088} \\
\noindent{Movements for Piano and Orchestra (1959), K094} \\

\section{Chamber}
\noindent{Pastorale for Violin, Clarinet, and Piano (1907), K004} \\
\noindent{Three Pieces for String Quartet (1914), K019} \\
\noindent{Suite Italienne for Violin and Piano (1932–33), K048} \\
\noindent{Duo Concertant for Violin and Piano (1932), K052} \\
\noindent{Septet (1953), K083} \\
\noindent{Elegy for Solo Viola (1944), K067} \\
\noindent{Concertino for Twelve Instruments (1952), K081} \\
\noindent{Double Canon for String Quartet (1959), K096} \\
\noindent{Septet for Mixed Instruments (1953), K083} \\
\noindent{Epitaphium for Flute, Clarinet, and Harp (1959), K095} \\
\noindent{Introitus for Solo Violin (1965), K106} \\

\section{Choral}
\noindent{Four Russian Peasant Songs for Women's Chorus and Four Horns (1914–17), K020} \\
\noindent{Les Noces (Svadebka), ballet-cantata for Soloists, Chorus, 4 Pianos, and Percussion (1914–17; rev. 1923), K024} \\
\noindent{Pater Noster, for Chorus (1926), K039} \\
\noindent{Symphony of Psalms, for Chorus and Orchestra (1930), K046} \textit{Tenebrae, BBC SO, Nigel Short}\marginnote{max}  \\
\noindent{Credo, for Chorus (1932), K053} \\
\noindent{Ave Maria, for Chorus (1934), K054} \\
\noindent{Cantata: Babel, for Narrator, Chorus, and Orchestra (1944), K068} \\
\noindent{Mass, for Mixed Chorus and Double Wind Quintet (1944–48), K077} \\
\noindent{Cantata: The King of the Stars (1948), K078} \\
\noindent{Threni: Lamentations of the Prophet Jeremiah, for Soloists, Chorus, and Orchestra (1957–58), K091} \\
\noindent{The Flood, Biblical Story with Narrator, Chorus, and Orchestra (1962), K099} \\
\noindent{Requiem Canticles, for Alto, Bass, Chorus, and Orchestra (1966), K109} \\

\section*{Vocal}
\noindent{Pastorale, for Voice and Piano (1907) K001} \\
\noindent{Two Poems of Paul Verlaine, for Voice and Piano (1907–08), K002} \\
\noindent{Pastorale, for Voice and Small Orchestra (1907, arr. 1933), K001a} \\
\noindent{Four Songs, for Voice and Piano (1906–09), K004} \\
\noindent{Two Poems of Konstantin Balmont, for Voice and Piano (1908), K006} \\
\noindent{Three Japanese Lyrics, for Voice and Chamber Orchestra (1912–13), K014} \\
\noindent{Pribaoutki, for Voice and 8 Instruments (1914), K019} \\
\noindent{Berceuses du chat (Cat’s Cradle Songs), for Voice and 3 Clarinets (1915–16), K021} \\
\noindent{Four Russian Songs, for Voice and Piano (1918–19), K030} \\
\noindent{Three Little Songs, for Voice and Piano (1919), K031} \\
\noindent{Two Balmont Songs, for Voice and Piano (1911), K015} \\
\noindent{Four Russian Peasant Songs, for Voice and Piano (1917), K020a} \\
\noindent{Four Russian Songs, for Voice and 8 Instruments (1919), K032} \\
\noindent{Three Songs from William Shakespeare, for Voice and Flute, Clarinet, and Viola (1953), K083} \\
\noindent{Elegy for J.F.K., for Baritone and 3 Clarinets (1964), K104} \\

\section*{Keyboard}
\noindent{Scherzo, for Piano (1902) K001} \\
\noindent{Piano Sonata in F-sharp minor (1903–04) K001} \\
\noindent{Four Etudes, for Piano (1908) K004} \\
\noindent{Piano Sonata in D minor (1907) K006} \\
\noindent{Eight Easy Pieces, for Piano (1914) K010} \\
\noindent{Piano-Rag-Music (1919) K035} \\
\noindent{Serenade en la, for Piano (1925) K040} \\
\noindent{Piano Sonata in F-sharp (1924) K048} \\
\noindent{Tango, for Piano (1940) K062} \\
\noindent{Sonata, for Two Pianos (1943–44) K072} \\
\noindent{Circus Polka, for Piano (1942) K066} \\





